\section{Related work}
\label{sec:related-work}

\subsection{Post-training quantization for detection}

Classical PTQ methods optimize scales, rounding, or local reconstruction to
reduce the gap between full-precision and quantized networks
\citep{nagel2020adaround,li2021brecq,hubara2021smallcalibration}. Object
detection is especially sensitive because classification and localization
heads, feature pyramids, and highly imbalanced foreground/background responses
must remain jointly calibrated. Fully Quantized Networks established low-bit
detection as a distinct problem\citep{li2019fullyquantized}; PD-Quant introduced
a prediction-difference objective\citep{liu2023pdquant}; and DetPTQ selects
layer-specific reconstruction metrics using an object-detection output loss
\citep{niu2023detptq}. More recently, InlierQ suppresses task-irrelevant
activation anomalies while retaining informative volumes with a small,
label-free calibration set\citep{kim2026inlierq}. These methods target clean
quantization fidelity. Our work addresses the complementary question of how to
evaluate quantized detectors when the input distribution is deliberately
corrupted.

\subsection{INT8, FP8, and executable treatments}

Integer-only inference provides an established hardware path for efficient
neural networks\citep{jacob2018integer}. FP8 formats trade mantissa precision
for exponent range and can preserve activation outliers that are difficult for
a fixed integer scale\citep{micikevicius2022fp8}. FP8 PTQ studies have reported
competitive accuracy and efficiency across model classes\citep{shen2024fp8},
while flexible-format work shows that the preferred 8-bit representation can
vary across layers and architectures\citep{zhang2023flexible8bit}. This
literature motivates a direct comparison but also cautions against reducing an
engine to its nominal datatype. Quantizer coverage, scale granularity,
calibration, graph transformations, and unsupported operations define the
actual treatment. We therefore inspect Q/DQ graph structure and interpret all
results as conditional on the recorded ModelOpt/TensorRT pipelines.

\subsection{Corruption robustness and quantized models}

Common-corruption benchmarks established standardized synthetic shifts for
image classification\citep{hendrycks2019corruptions}. Detector studies then
showed that localization systems can fail differently from classifiers under
weather, noise, blur, and camera-motivated distortions
\citep{michaelis2019detectorrobustness,liu2024realworldcorruptions}. Work in
\emph{Image and Vision Computing} has also addressed adverse-weather detection,
test-time adaptation, compact edge localization, and small-object failure modes
\citep{luo2024lalm,zhang2025ddt,su2024yolic,zheng2024smallobjectreview}.
Small objects are particularly vulnerable because fewer pixels support class and
box estimates\citep{cheng2023small}, but low absolute small-object AP does not
by itself imply a precision-specific small-object interaction.

Research directly joining quantization and distribution shift is emerging.
Karimov et al. compare detector quantization under input degradations
\citep{karimov2025robustness}; TTAQ adapts PTQ under continuous domain change
\citep{xiao2024ttaq}; and Recti-Q targets out-of-distribution robustness through
feature-space rectification\citep{araghi2026rectiq}. Related adversarial work
finds that apparently inconsistent robustness conclusions can arise from
changes in the quantization pipeline itself\citep{li2024quantrobustness}. These
studies reinforce the need to state the complete treatment and the exact
contrast being estimated.

\subsection{Positioning of the present study}

Our work does not propose a new quantizer, corruption generator, or detector.
It contributes a measurement protocol for executable engines. Three design
choices distinguish it from a raw corrupted-accuracy comparison. First, the
clean and corrupted inputs are codec-matched and shared byte-for-byte across
formats. Second, the clean FP8--INT8 gap is removed inside a common image
resample, yielding a direct difference in differences rather than a subtraction
of separately summarized results. Third, the signed interaction is interpreted
only after matched-clean fidelity and absolute corrupted AP have been checked.
The resulting protocol is designed to be reusable across architectures and
runtimes, including cases in which a PTQ recipe fails to transfer.
